\documentclass[a4paper,12pt]{article}
\usepackage[utf8]{inputenc}
\usepackage[russian]{babel}
\usepackage{amsmath, amssymb}
\usepackage{geometry}
\geometry{top=2cm, bottom=2cm, left=2cm, right=2cm}

\title{Задание 3.1}
\author{}
\date{}

\begin{document}
	
	\maketitle
	
	\section*{Вывод формулы (12)}
	
	Методом неопределённых коэффициентов найдём формулу для \( f''(x) \) по значениям 
	\( f(x), f(x+h), f(x+2h), f(x+3h) \):
	
	\[
	f''(x) \approx \alpha f(x) + \beta f(x+h) + \gamma f(x+2h) + \delta f(x+3h).
	\]
	
	\subsection*{Разложение в ряд Тейлора в точке \( x \)}
	\[
	\begin{aligned}
		f(x) &= f, \\
		f(x+h) &= f + hf' + \frac{h^2}{2}f'' + \frac{h^3}{6}f''' + \frac{h^4}{24}f^{(4)} + O(h^5), \\
		f(x+2h) &= f + 2hf' + 2h^2f'' + \frac{4h^3}{3}f''' + \frac{2h^4}{3}f^{(4)} + O(h^5), \\
		f(x+3h) &= f + 3hf' + \frac{9h^2}{2}f'' + \frac{9h^3}{2}f''' + \frac{27h^4}{8}f^{(4)} + O(h^5).
	\end{aligned}
	\]
	
	\subsection*{Линейная комбинация}
	\[
	\alpha f(x) + \beta f(x+h) + \gamma f(x+2h) + \delta f(x+3h) =
	\]
	\[
	= (\alpha + \beta + \gamma + \delta)f +
	h(\beta + 2\gamma + 3\delta)f' +
	\frac{h^2}{2}(\beta + 4\gamma + 9\delta)f'' +
	\]
	\[
	+ \frac{h^3}{6}(\beta + 8\gamma + 27\delta)f''' +
	\frac{h^4}{24}(\beta + 16\gamma + 81\delta)f^{(4)} + O(h^5).
	\]
	
	\subsection*{Система уравнений}
	\[
	\begin{cases}
		\alpha + \beta + \gamma + \delta = 0 & (1) \\[4pt]
		\beta + 2\gamma + 3\delta = 0 & (2) \\[4pt]
		\dfrac{1}{2}(\beta + 4\gamma + 9\delta) = 1 & (3) \\[4pt]
		\beta + 8\gamma + 27\delta = 0 & (4)
	\end{cases}
	\]
	
	\subsection*{Решение}
	Из (2): \( \beta = -2\gamma - 3\delta \).
	
	Подставим в (3):
	\[
	\frac{1}{2}\bigl((-2\gamma-3\delta) + 4\gamma + 9\delta\bigr) = 1
	\;\Rightarrow\;
	\frac{1}{2}(2\gamma + 6\delta) = 1
	\;\Rightarrow\;
	\gamma + 3\delta = 1. \quad (3')
	\]
	
	Подставим в (4):
	\[
	(-2\gamma-3\delta) + 8\gamma + 27\delta = 0
	\;\Rightarrow\;
	6\gamma + 24\delta = 0
	\;\Rightarrow\;
	\gamma + 4\delta = 0. \quad (4')
	\]
	
	Вычтем (4') из (3'):
	\[
	(\gamma + 3\delta) - (\gamma + 4\delta) = 1 - 0
	\;\Rightarrow\;
	-\delta = 1
	\;\Rightarrow\;
	\delta = -1.
	\]
	
	Из (4'): \( \gamma + 4(-1) = 0 \) \( \Rightarrow \) \( \gamma = 4 \).
	
	Из (2): \( \beta = -2\cdot4 - 3\cdot(-1) = -8 + 3 = -5 \).
	
	Из (1): \( \alpha - 5 + 4 - 1 = 0 \) \( \Rightarrow \) \( \alpha = 2 \).
	
	\[
	\boxed{\alpha = 2,\quad \beta = -5,\quad \gamma = 4,\quad \delta = -1}
	\]
	
	\subsection*{Формула (12) и её погрешность}
	\[
	\boxed{f''(x) = \frac{2f(x) - 5f(x+h) + 4f(x+2h) - f(x+3h)}{h^2} + \frac{11}{12}h^2 f^{(4)}(\xi)},\quad \xi \in (x, x+3h).
	\]
	
	\section*{Вывод формулы (13)}
	
	Замена \( h \to -h \) в (12):
	
	\[
	\boxed{f''(x) = \frac{2f(x) - 5f(x-h) + 4f(x-2h) - f(x-3h)}{h^2} + \frac{11}{12}h^2 f^{(4)}(\xi)},\quad \xi \in (x-3h, x).
	\]
	
\end{document}